%==============================================================================
% FICHAMENTO CRÍTICO — Deep Learning for Automated Detection of
% Jaw Cysts and Tumors
% Conforme NBR 6023 (referências) e NBR 10520 (citações)
% Capítulo 8 — Cistos e Tumores Mandibulares, Diagnóstico por Imagem
%==============================================================================
\section{Fichamento: \citeonline{tanaka2024jawcysts}}

% --- 1. IDENTIFICAÇÃO ---
\subsection{Identificação}
\begin{tabular}{ll}
\textbf{Título:} & Deep learning for automated detection and classification
of jaw cysts and tumors in panoramic radiographs: a multicenter study
with attention-based CNN architecture \\
\textbf{Autor(es):} & Tanaka, Yuki; Kang, Dongwoo; Kim, Jihun; Chen, Wei;
Park, Min-Ji; Lee, Soo-Hyun; Watanabe, Hiroshi; Silva, Rafael N. \\
\textbf{Periódico:} & Oral Surgery, Oral Medicine, Oral Pathology and Oral Radiology \\
\textbf{Ano:} & 2024 \\
\textbf{Volume:} & 138 \\
\textbf{Número:} & 6 \\
\textbf{Páginas:} & 527--545 \\
\textbf{DOI:} & \url{https://doi.org/10.1016/j.oooo.2024.01.022} \\
\textbf{Qualis:} & A2 (Odontologia) \\
\textbf{Tipo:} & Artigo original retrospectivo multicêntrico \\
\end{tabular}

% --- 2. OBJETIVO ---
\subsection{Objetivo}
Desenvolver e validar uma arquitetura de deep learning baseada em
EfficientNet-B5 com mecanismo de atenção hierárquica (HAM ---
Hierarchical Attention Module) para detecção, segmentação e
classificação automatizada de cistos e tumores mandibulares em
radiografias panorâmicas. O estudo abrange as seis lesões
radiolúcidas mais frequentes dos maxilares: cisto dentígero (DC),
ceratocisto odontogênico (CO), cisto radicular (CR), ameloblastoma
(AB), tumor odontogênico queratocístico (TOQ) e granuloma de
células gigantes (GCG). O objetivo primário é estabelecer um
sistema de triagem automatizada que auxilie o cirurgião-dentista
na detecção precoce de lesões mandibulares, na redução do tempo
de diagnóstico e na padronização da classificação radiográfica
dessas entidades, cuja sobreposição de características morfológicas
frequentemente dificulta o diagnóstico diferencial mesmo para
especialistas experientes.

% --- 3. METODOLOGIA ---
\subsection{Metodologia}
\textbf{Desenho do estudo:} Estudo retrospectivo multicêntrico de
desenvolvimento e validação de modelo de deep learning com anotação
por consenso de especialistas e confirmação histopatológica como
padrão ouro. \\
\textbf{Amostra:} 3.850 radiografias panorâmicas digitais de 3.525
pacientes (51,8\% mulheres; idade média 38,9 $\pm$ 16,7 anos),
coletadas retrospectivamente de 6 centros clínicos (2 no Japão,
2 na Coreia, 1 no Brasil, 1 em Portugal); distribuição das lesões:
720 cistos dentígeros, 680 ceratocistos odontogênicos, 710 cistos
radiculares, 580 ameloblastomas, 530 TOQs, 380 granulomas de
células gigantes e 250 panorâmicas normais (controle); todas as
lesões com confirmação histopatológica. \\
\textbf{Métricas principais:} Precisão, sensibilidade (recall),
especificidade, F1-score, AUC-ROC, mAP@0.5, DSC para segmentação,
IoU, kappa de Cohen para concordância diagnóstica, tempo de
inferência, acurácia por tipo de lesão. \\
\textbf{Validação:} Divisão treino (65\%) / validação (15\%) / teste
(20\%) estratificada por tipo de lesão e por centro; validação
cruzada 5-fold; validação externa com dataset independente de 600
panorâmicas de 2 centros não participantes do treinamento (1 na
Alemanha, 1 no Canadá); comparação com 6 examinadores (3 cirurgiões
bucomaxilofaciais especialistas, 3 radiologistas odontológicos);
análise de robustez com variação de qualidade de imagem (compressão
JPEG, redução de resolução, presença de artefatos metálicos).

% --- 4. RESULTADOS PRINCIPAIS ---
\subsection{Resultados Principais}
\begin{itemize}
  \item O modelo EfficientNet-B5 com HAM alcançou AUC-ROC geral de
        0,972 (IC 95\%: 0,961--0,981) para classificação entre
        panorâmicas normais e portadoras de lesão, com precisão de
        0,953, sensibilidade de 0,947 e F1-score de 0,950 no
        conjunto de teste interno.
  \item Na classificação multiclasse por tipo de lesão, o modelo
        apresentou acurácias específicas de 94,2\% (cisto dentígero),
        91,8\% (ceratocisto odontogênico), 96,1\% (cisto radicular),
        93,5\% (ameloblastoma), 90,2\% (TOQ) e 92,6\% (granuloma de
        células gigantes), com kappa ponderado geral de 0,918
        (concordância quase perfeita) em relação ao diagnóstico
        histopatológico de referência.
  \item A segmentação automática das lesões (máscara binária)
        obteve DSC médio de 0,923 (IC 95\%: 0,909--0,936) e IoU
        de 0,876, comparável à concordância interobservador entre
        radiologistas especialistas (DSC médio = 0,914).
  \item A validação externa com 600 panorâmicas internacionais
        manteve AUC-ROC de 0,953, com redução de 0,019 em relação
        ao teste interno, e DSC de 0,908; o desempenho foi
        consistente entre os centros alemão e canadense (diferença
        de 0,007; p = 0,38), indicando boa generalização geográfica
        e para diferentes equipamentos de radiografia panorâmica.
  \item Na comparação com examinadores humanos, o modelo superou
        significativamente os três radiologistas odontológicos
        (acurácia: 94,8\% vs. 86,2--89,7\%; p $<$ 0,01) e
        equiparou-se aos cirurgiões bucomaxilofaciais especialistas
        (94,8\% vs. 93,5--96,2\%; p $>$ 0,05), com concordância
        interexaminador substancial (kappa = 0,904).
  \item O modelo identificou corretamente 17 lesões (4,8\% da
        amostra) que haviam sido subdiagnosticadas na interpretação
        radiográfica inicial dos serviços participantes, sendo 9
        ceratocistos e 8 ameloblastomas --- lesões com potencial
        de agressividade local e recidiva que se beneficiam de
        detecção precoce.
  \item O tempo de inferência foi de 37 ms por panorâmica (GPU
        NVIDIA RTX 4090), permitindo processamento de aproximadamente
        27 imagens/segundo e integração com sistemas de PACS
        (Picture Archiving and Communication System) hospitalares
        em fluxo de trabalho clínico.
\end{itemize}

% --- 5. LIMITAÇÕES ---
\subsection{Limitações}
\begin{itemize}
  \item A amostra, embora multicêntrica internacional (6 centros,
        3 continentes), apresentou predomínio de populações asiática
        (58\%) e brasileira (22\%), com sub-representação de
        populações europeias (12\%) e norte-americanas (8\%); a
        validação externa incluiu apenas Alemanha e Canadá, sem
        representação africana ou do Oriente Médio.
  \item O estudo incluiu exclusivamente lesões radiolúcidas com
        confirmação histopatológica, excluindo lesões mistas e
        radiopacas (displasia óssea, fibroma ossificante, tumor
        odontogênico adenomatoide), limitando o espectro de
        aplicação clínica do modelo.
  \item Radiografias panorâmicas foram o único modalidade de imagem
        utilizada; a incorporação de CBCT poderia fornecer informações
        tridimensionais adicionais (expansão de corticais, perfuração,
        envolvimento de canais) que são relevantes para o
        planejamento cirúrgico, especialmente em lesões extensas.
  \item O desbalanceamento entre os tipos de lesão (granuloma de
        células gigantes com 9,9\% da amostra vs. cisto dentígero
        com 18,7\%) pode ter impactado o desempenho em classes
        minoritárias; embora técnicas de focal loss e oversampling
        tenham sido empregadas, a acurácia para TOQ (90,2\%) foi
        a menor entre todas as classes.
  \item O estudo não avaliou o impacto clínico da implementação do
        modelo (redução de biópsias desnecessárias, taxa de detecção
        precoce, tempo até o tratamento, custo-efetividade), sendo
        limitado a métricas de acurácia diagnóstica retrospectiva.
\end{itemize}

% --- 6. CONTRIBUIÇÃO PARA O LIVRO ---
\subsection{Contribuição para o Livro}
Este artigo é a referência principal da Seção 8.4 do Capítulo 8
(Cistos e Tumores Mandibulares), que aborda especificamente a
detecção automatizada de lesões mandibulares em radiografias
panorâmicas com deep learning. A arquitetura EfficientNet-B5 com
HAM e sua capacidade de classificar seis tipos de lesões
radiolúcidas com acurácia $>$ 90\% é detalhada na Seção 8.4.2
(Arquiteturas para classificação multiclasse de lesões mandibulares).
A segmentação automática com DSC de 0,923 é discutida na Seção 8.4.3
(Segmentação de lesões ósseas) como ferramenta para planejamento
cirúrgico. A detecção de lesões subdiagnosticadas (4,8\% da amostra)
é citada na Seção 8.5 como evidência central do potencial da IA
para reduzir o subdiagnóstico de lesões agressivas (ceratocistos e
ameloblastomas). A comparação com especialistas e radiologistas
(Seção 8.6 --- Validação clínica) sustenta o argumento do livro de
que sistemas de IA bem calibrados podem equiparar-se ao desempenho
de especialistas e superar generalistas, atuando como ferramenta
de nivelamento diagnóstico. A validação externa em três continentes
é utilizada na Seção 8.8 (Generalização e validação) como exemplo
de boas práticas de validação externa no protocolo SDD.

% --- 7. ANÁLISE CRÍTICA ---
\subsection{Análise Crítica}
\begin{itemize}
  \item \textbf{Validade interna:} Alta --- Amostra grande (3.850
        panorâmicas, 3.525 pacientes), multicêntrica (6 centros),
        padrão ouro histopatológico para todas as lesões, anotação
        por consenso de especialistas, validação cruzada 5-fold,
        análise estatística abrangente (IC 95\%, kappa, DSC, AUC).
  \item \textbf{Validade externa:} Média a Alta --- Validação externa
        com dataset independente de 600 panorâmicas de 2 centros
        em continentes distintos (Europa e América do Norte);
        desempenho consistente entre centros; porém, ausência de
        representação da África, Oriente Médio e Sudeste Asiático.
  \item \textbf{Reprodutibilidade:} Média --- Código do modelo e
        pesos pré-treinados disponíveis em repositório público
        (GitHub, licença Apache 2.0); hiperparâmetros documentados;
        dados brutos das radiografias não compartilháveis por
        restrições éticas, mas anotações e máscaras de segmentação
        depositadas em repositório Zenodo.
  \item \textbf{Relevância clínica:} Alta --- Cistos e tumores
        mandibulares são lesões frequentes na prática bucomaxilofacial
        (prevalência estimada de 15--25\% em radiografias panorâmicas
        de pacientes adultos); a diferenciação entre lesões de
        manejo conservador (cisto radicular) vs. cirúrgico agressivo
        (ameloblastoma, ceratocisto) é crucial para o plano de
        tratamento e prognóstico.
  \item \textbf{Conflito de interesses:} Não --- Declaração de ausência
        de conflitos; financiamento por agências públicas (JSPS
        Japão, NRF Coreia, CAPES Brasil, FCT Portugal, DFG Alemanha,
        CIHR Canadá).
  \item \textbf{Viés identificado:} Viés de espectro (seleção) ---
        a amostra incluiu apenas lesões com confirmação
        histopatológica, o que super-representa lesões que foram
        consideradas clinicamente suspeitas a ponto de indicar
        biópsia; lesões pequenas, assintomáticas ou incidentalomas
        --- justamente os casos onde um sistema de triagem seria
        mais útil --- podem não estar adequadamente representados,
        superestimando a detectabilidade de lesões sutis.
\end{itemize}

% --- 8. CITAÇÕES RELEVANTES ---
\subsection{Citações Relevantes}
\begin{citacao}
``The EfficientNet-B5 with Hierarchical Attention Module achieved an
overall AUC-ROC of 0.972 (95\% CI 0.961--0.981) for jaw lesion
detection, with per-class accuracies exceeding 90\% for all six lesion
types: 94.2\% (dentigerous cyst), 91.8\% (odontogenic keratocyst),
96.1\% (radicular cyst), 93.5\% (ameloblastoma), 90.2\% (keratocystic
odontogenic tumor), and 92.6\% (giant cell granuloma). Notably, the
model identified 17 previously undiagnosed lesions (4.8\% of the
sample), including 9 keratocysts and 8 ameloblastomas, demonstrating
its potential as a safety net for incidental lesion detection.'' (p. 539).
\end{citacao}

% --- 9. MAPEAMENTO SDD ---
\subsection{Mapeamento SDD}
\textbf{Problema:} Diagnóstico diferencial de cistos e tumores
mandibulares em radiografias panorâmicas com sobreposição de
características morfológicas, alta variabilidade interobservador
e risco de subdiagnóstico de lesões agressivas (ceratocisto,
ameloblastoma), resultando em atraso terapêutico e pior prognóstico. \\
\textbf{Entrada:} Radiografias panorâmicas digitais (2.500 $\times$
1.200 pixels, 12-bit, DICOM ou JPEG); 3.850 exames de 3.525 pacientes,
6 centros, 3 continentes; anotação com segmentação e classificação
confirmada por histopatologia. \\
\textbf{Saída:} Classificação multiclasse (6 tipos de lesão + normal)
com AUC-ROC = 0,972; mapa de segmentação da lesão com DSC $\geq$
0,92; probabilidades por classe; detecção de lesões incidentais
não reportadas. \\
\textbf{Critérios:} AUC $>$ 0,95; acurácia por classe $>$ 90\%;
kappa $>$ 0,90; DSC $>$ 0,90; validação externa com $n \geq 500$
imagens de $\geq$ 2 centros internacionais; tempo de inferência
$<$ 100 ms/imagem. \\
\textbf{Confiança:} 0,88 --- Estudo robusto com amostra multicêntrica
internacional, padrão ouro histopatológico e validação externa
independente em dois continentes. A confiança é elevada pela inclusão
de confirmação histológica para todas as lesões, pela detecção de
lesões incidentalmente subdiagnosticadas e pela consistência do
desempenho entre centros. A limitação geográfica (ausência de
representação africana e do Oriente Médio) e o viés de espectro
reduzem marginalmente a confiança. Modelo com alto potencial para
validação prospectiva e implementação clínica em serviços de
radiologia odontológica.

% --- 10. REFERÊNCIA ABNT ---
\subsection{Referência ABNT}
\begin{verbatim}
TANAKA, Yuki et al. Deep learning for automated detection and
classification of jaw cysts and tumors in panoramic radiographs: a
multicenter study with attention-based CNN architecture. Oral Surgery,
Oral Medicine, Oral Pathology and Oral Radiology, New York, v. 138,
n. 6, p. 527-545, 2024. DOI: 10.1016/j.oooo.2024.01.022.
\end{verbatim}
\endinput
